% ============================================================
%  journal_paper.tex
%  Automated Blood Group Detection Using Computer Vision and
%  Transfer Learning on Mobile-Optimized Deep Neural Networks
%
%  IEEE Transactions / IEEE Access Journal Style
%  Compile with:  pdflatex journal_paper.tex  (run twice)
% ============================================================
\documentclass[journal]{IEEEtran}

% ---- Packages -----------------------------------------------
\usepackage[T1]{fontenc}
\usepackage[utf8]{inputenc}
\usepackage{amsmath,amssymb}
\usepackage{graphicx}
\usepackage{booktabs}
\usepackage{array}
\usepackage{url}
\usepackage{cite}

% =============================================================
\begin{document}

% ---- Title & Author -----------------------------------------
\title{Automated Blood Group Detection Using Computer Vision and
       Transfer Learning on Mobile-Optimized Deep Neural Networks:
       An End-to-End Pipeline for Point-of-Care Screening}

\author{\IEEEauthorblockN{Sayon Mitra}
\IEEEauthorblockA{Independent Researcher, AI and Biomedical Engineering\\
Email: sayon@sayonedu.in}}

\maketitle

% =============================================================
\begin{abstract}
Accurate and rapid blood group identification is a critical
prerequisite in transfusion medicine, emergency care, and
perioperative management.
Conventional haemagglutination-based typing is time-consuming,
operator-dependent, and unsuitable for point-of-care use in
resource-limited settings.
This paper presents a comprehensive, end-to-end AI-driven blood
group detection system that combines COCO-format bounding-box
region extraction, MobileNetV2 transfer learning, and deterministic
decision logic to classify ABO and Rh(D) blood groups directly from
ambient-light photographs of standard disposable blood-typing cards.
The pipeline autonomously extracts three reagent zones (Anti-A,
Anti-B, Anti-D), classifies agglutination intensity in each zone
as Strong Clumping, Medium Clumping, or No Clumping, and applies
deterministic ABO\,+\,Rh inference to produce the final blood group
designation.
A curated dataset of 7,836 annotated card images (6,456 training /
997 validation / 383 test) was assembled via Roboflow and
encompasses all four ABO types with both Rh polarities.
The system achieves $\approx$94--95\,\% zone-level classification
accuracy and $\approx$91--93\,\% end-to-end ABO+Rh accuracy, with
a macro-average F1-score of 0.947 on 1,149 zone-level test crops.
The 14\,MB MobileNetV2 backbone supports deployment on
resource-constrained mobile devices via TensorFlow Lite, enabling
sub-50\,ms inference on mid-range Android hardware.
A systematic error analysis reveals that the primary failure mode
is weak-D (D\textsuperscript{u}) phenotype misclassification,
arising from physiologically faint agglutination in the Anti-D zone.
This work also forms the AI inference core of \emph{Gyanbondhu}, an
AI-powered biodegradable blood group testing kit designed for
autonomous field deployment in low-resource healthcare settings.
This system is intended for assistive screening only and is not
a substitute for laboratory-validated clinical diagnosis.
\end{abstract}

\begin{IEEEkeywords}
blood group detection, transfer learning, MobileNetV2,
COCO annotation, computer vision, haemagglutination,
point-of-care diagnostics, AI in healthcare, Rh factor,
biodegradable testing kit
\end{IEEEkeywords}

% =============================================================
\section{Introduction}
\label{sec:intro}

\IEEEPARstart{B}{lood} transfusion incompatibility is one of the
leading preventable causes of transfusion-related fatalities worldwide.
The ABO and Rh(D) blood group systems are the two most clinically
significant antigen systems; mismatching in either can trigger
life-threatening acute or delayed haemolytic transfusion
reactions~\cite{lippi2011}.
Despite advances in pre-transfusion testing, manual haemagglutination
typing remains susceptible to inter-observer variability of 2--5\,\%
due to subjective interpretation of agglutination
strength~\cite{harmening2019}.

The challenge is acute in resource-limited settings.
Rural healthcare centres, military field hospitals, disaster-response
units, and emergency pre-hospital services frequently lack the
laboratory infrastructure, trained staff, and reagent supply chains
required for conventional blood typing.
Delayed or erroneous blood group identification in these contexts
can be life-threatening, particularly in mass casualty events
or emergency surgeries.

The rapid proliferation of high-resolution smartphone cameras
and the availability of lightweight deep learning frameworks
(TensorFlow Lite, ONNX Runtime, PyTorch Mobile) create a compelling
opportunity to deploy AI-powered blood group typing at the point of
care without requiring dedicated laboratory equipment.
A system capable of classifying blood group from a photograph of a
standard disposable typing card taken under ordinary ambient light
would substantially reduce the diagnostic bottleneck in field
and low-resource settings.

This paper presents an end-to-end pipeline for automated blood group
detection that (i) extracts reagent zones from full blood-card images
using COCO bounding-box annotations, (ii) classifies agglutination
intensity in each zone using MobileNetV2 transfer learning, and
(iii) applies deterministic logic to derive the ABO+Rh blood group.
The system is trained on a large, curated, Roboflow-annotated dataset
and evaluated rigorously on held-out test data.

The work also provides the deep-learning inference engine for
\emph{Gyanbondhu}---an AI-powered biodegradable blood group testing
kit that integrates the camera-based classifier with a low-cost
disposable card and a mobile application for autonomous, field-deployable
blood typing without laboratory infrastructure.

\subsection{Contributions}

The specific contributions of this work are:
\begin{enumerate}
  \item \textbf{End-to-end automated pipeline.}
        A complete card-to-blood-group system: COCO region extraction
        $\rightarrow$ MobileNetV2 clumping classification
        $\rightarrow$ deterministic ABO\,+\,Rh decision logic,
        operating on unconstrained ambient-light photographs.
  \item \textbf{Large-scale annotated dataset.}
        A curated, Roboflow-annotated dataset of 7,836 blood-card
        images (6,456 training / 997 validation / 383 test) covering
        all ABO types and both Rh polarities---the largest publicly
        described dataset for this task.
  \item \textbf{Mobile-deployable model.}
        A MobileNetV2-based model (14\,MB) with a documented
        TensorFlow Lite conversion roadmap for $<$50\,ms on-device
        inference on mid-range Android devices.
  \item \textbf{Error analysis and failure characterisation.}
        A systematic analysis of error modes, with identification
        of weak-D phenotype misclassification as the primary
        failure source and a confidence-threshold mitigation strategy.
  \item \textbf{Real-world deployment context.}
        Integration architecture for \emph{Gyanbondhu}, a
        biodegradable blood group testing kit targeting low-cost
        healthcare in rural and field settings.
\end{enumerate}

% =============================================================
\section{Related Work}
\label{sec:related}

\subsection{Traditional Blood Group Detection Methods}

Haemagglutination tests, in which red blood cells mixed with
specific antisera clump (agglutinate) when the corresponding antigen
is present, have constituted the clinical gold standard for blood
group typing since Landsteiner's discovery of the ABO system in
1901~\cite{landsteiner1901}.
Three standard formats exist: glass-slide agglutination
(qualitative, rapid, low-cost), test-tube agglutination (more
sensitive, semi-quantitative), and microplate (gel card or column
agglutination) techniques (highly sensitive, amenable to automation).
All formats share a common limitation: reaction-strength
interpretation is subjective and operator-dependent, with reported
inter-observer variability of 2--5\,\%~\cite{harmening2019}.
Weak or partial agglutination is particularly susceptible to
misclassification, as is weak-D (D\textsuperscript{u}), a
phenotypic variant of Rh(D) that produces faint agglutination
with standard anti-D reagents~\cite{harmening2019}.

\subsection{Computer-Vision Approaches}

Early automated systems relied on classical image-processing
heuristics.
Saif~\emph{et al.}~\cite{saif2016} used HSV colour thresholding
and blob detection to segment agglutinated regions on glass slides,
achieving 88\,\% accuracy on a 200-image dataset under controlled
illumination.
Khan~\emph{et al.}~\cite{khan2018} employed a support-vector machine
(SVM) classifier on hand-crafted texture features derived from
grey-level co-occurrence matrices (GLCM), reporting 91.3\,\% accuracy
on a 450-image controlled dataset.
Both approaches are sensitive to illumination change and require
single-reagent pre-cropped images, limiting real-world applicability.

\subsection{Deep Learning Approaches}

Convolutional neural networks substantially improved accuracy.
Yildiz and Aslan~\cite{yildiz2020} applied ResNet-50 transfer
learning to pre-cropped single-reagent images, achieving 93.7\,\%
on a 1,500-image dataset under controlled laboratory conditions.
Ahmed~\emph{et al.}~\cite{ahmed2021} used VGG-16 with data
augmentation on a 1,200-image dataset, reporting 94.2\,\% accuracy.
Despite strong performance, both studies used large, labour-intensive
CNN backbones (25--138\,M parameters) unsuitable for mobile deployment,
and neither addressed the problem of extracting reagent regions from
full, uncropped card images.

\subsection{Comparison and Gap Analysis}

Table~\ref{tab:comparison} summarises prior work alongside the
proposed system.
The proposed system is distinguished by: (i) a complete card-level
detection pipeline that requires no manual pre-cropping; (ii) a
mobile-friendly backbone (3.4\,M parameters) enabling on-device
inference; (iii) a dataset nearly five times larger than any prior
study; and (iv) systematic error analysis addressing failure cases.
No prior published work has reported an end-to-end automated pipeline
from raw card photograph to final blood group designation under
ambient light conditions.

\begin{table}[htbp]
  \caption{Comparison with Prior Blood Group Detection Systems}
  \label{tab:comparison}
  \centering
  \renewcommand{\arraystretch}{1.1}
  \begin{tabular}{lp{1.4cm}p{0.8cm}p{1.0cm}p{1.0cm}}
    \toprule
    \textbf{System} & \textbf{Method} & \textbf{Acc.} &
    \textbf{Dataset} & \textbf{Full Pipeline} \\
    \midrule
    Saif~\emph{et al.}~\cite{saif2016}    & HSV+Blob   & 88.0\,\% & 200    & No \\
    Khan~\emph{et al.}~\cite{khan2018}    & SVM+GLCM   & 91.3\,\% & 450    & No \\
    Yildiz~\emph{et al.}~\cite{yildiz2020}& ResNet-50  & 93.7\,\% & 1,500  & No \\
    Ahmed~\emph{et al.}~\cite{ahmed2021}  & VGG-16     & 94.2\,\% & 1,200  & No \\
    \textbf{Proposed}                     & \textbf{MNv2} & \textbf{94--95\,\%} & \textbf{7,836} & \textbf{Yes} \\
    \bottomrule
  \end{tabular}
\end{table}

% =============================================================
\section{Methodology}
\label{sec:method}

\subsection{Dataset Collection and Annotation}

Images of commercially available disposable blood-typing cards were
collected under diverse ambient-light conditions (natural daylight,
fluorescent office lighting, and artificial warm light) and with
multiple card brands to maximise distributional coverage.
Each full-card image was annotated using the Roboflow platform in
COCO JSON format~\cite{lin2014coco}, with three rectangular bounding
boxes per image delineating the Anti-A, Anti-B, and Anti-D reagent
zones.
Ground-truth labels were assigned by visual inspection of each
zone crop by the researcher, with the following classification
criteria:
\begin{itemize}
  \item \textbf{Strong Clumping (class 0):} Dense, clearly positive
        agglutination with large visible aggregates and no free
        red-cell suspension.
  \item \textbf{Medium Clumping (class 1):} Moderate positive
        agglutination with visible but smaller aggregates and
        partial free-cell suspension.
  \item \textbf{No Clumping (class 2):} Negative reaction; uniform
        red-cell suspension with no visible aggregates.
\end{itemize}

The dataset statistics are given in Table~\ref{tab:dataset}.
The train/validation/test split ratio is approximately 82\,\%\,/\,13\,\%\,/\,5\,\%
by image count, yielding 19,368 training zone crops, 2,991
validation zone crops, and 1,149 test zone crops.
The split was stratified by blood group class to ensure representative
coverage in all partitions.

\begin{table}[htbp]
  \caption{Dataset Statistics}
  \label{tab:dataset}
  \centering
  \renewcommand{\arraystretch}{1.1}
  \begin{tabular}{lccc}
    \toprule
    \textbf{Split} & \textbf{Images} & \textbf{Zone Crops} & \textbf{Fraction} \\
    \midrule
    Train      & 6,456 & 19,368 & 82.4\,\% \\
    Validation & 997   & 2,991  & 12.7\,\% \\
    Test       & 383   & 1,149  & ~4.9\,\% \\
    \midrule
    \textbf{Total} & \textbf{7,836} & \textbf{23,508} & 100\,\% \\
    \bottomrule
  \end{tabular}
\end{table}

\subsection{Data Augmentation}

To improve generalisation across the expected range of real-world
imaging conditions, the training set was augmented online using
TensorFlow data pipeline operations:
random horizontal and vertical flip, brightness jitter ($\pm$0.15),
contrast jitter ($\pm$0.10), rotation ($\pm$15$^{\circ}$),
and CLAHE (Contrast-Limited Adaptive Histogram Equalisation)
normalisation.
Validation and test sets were preprocessed identically (resize to
$224\times224$, RGB normalisation to $[0, 1]$) but without
augmentation to ensure unbiased evaluation.
CLAHE was applied to reduce sensitivity to uneven illumination,
which is the most common source of image quality variation in
field conditions.

\subsection{Region Extraction Pipeline}

Given a full blood-card image, the corresponding COCO annotation
file is parsed to retrieve three bounding boxes for the Anti-A,
Anti-B, and Anti-D zones.
Each bounding box is applied to the image using OpenCV to extract
the reagent region crop.
Crops are resized to $224\times224$ pixels, converted from BGR to
RGB colour order, and normalised to the $[0.0, 1.0]$ floating-point
range before being passed to the classifier.
The extracted crops are organised into a TensorFlow-compatible
directory structure (one subdirectory per class) using
\texttt{image\_dataset\_from\_directory}, supporting efficient
batched loading during training.

An automatic zone detector (e.g., YOLOv8-nano or MobileNet-SSD)
is identified as a high-priority future replacement for the
current annotation-dependent extraction step, which currently
requires a COCO JSON file to be pre-prepared for each card image.

\subsection{Model Architecture}

Transfer learning is applied using MobileNetV2~\cite{sandler2018}
pretrained on ImageNet-1K as the feature extraction backbone.
MobileNetV2 employs inverted residual bottleneck blocks with
depthwise separable convolutions and linear bottlenecks, achieving
an efficient trade-off between accuracy and parameter count
(3.4\,M parameters, 14\,MB float32 weights)~\cite{sandler2018}.
This compares favourably with ResNet-50 (25.6\,M, $\approx$100\,MB)
and VGG-16 (138\,M, $\approx$528\,MB), both of which are unsuitable
for mobile deployment.

The MobileNetV2 backbone (154 layers, 2.3\,M trainable parameters
in the convolutional base) is frozen during the initial training
phase to preserve ImageNet representations.
A task-specific classification head is appended, consisting of:
GlobalAveragePooling2D (1280-dimensional output), Dropout(0.3),
Dense(128, ReLU activation), and Dense(3, Softmax activation).
Total trainable parameters during the initial phase:
$\approx$166\,K (head only).

An optional fine-tuning phase unfreezes the last 30 convolutional
layers of the MobileNetV2 backbone (excluding Batch-Normalisation
layers to preserve running statistics) and trains with a reduced
learning rate of $\eta{=}10^{-5}$ for domain adaptation to the
blood-card imaging domain.

\subsection{Training Configuration}

The model is compiled with the Adam optimiser, sparse categorical
cross-entropy loss, and a batch size of 32.
Table~\ref{tab:training} gives the full hyperparameter configuration.
Early stopping with patience 3 on validation accuracy prevents
overfitting; a learning-rate reduction-on-plateau callback
(factor 0.5, patience 2, minimum $10^{-6}$) is applied to
stabilise convergence.
Training converges within 5--8 epochs on the frozen backbone phase.

\begin{table}[htbp]
  \caption{Training Hyperparameters}
  \label{tab:training}
  \centering
  \renewcommand{\arraystretch}{1.1}
  \begin{tabular}{ll}
    \toprule
    \textbf{Hyperparameter} & \textbf{Value} \\
    \midrule
    Backbone                & MobileNetV2 (ImageNet pre-trained) \\
    Optimiser               & Adam \\
    Initial learning rate   & $1 \times 10^{-4}$ (phase 1); $1 \times 10^{-5}$ (fine-tune) \\
    Loss function           & Sparse Categorical Cross-Entropy \\
    Batch size              & 32 \\
    Maximum epochs          & 10 \\
    Early-stopping patience & 3 (val\_accuracy) \\
    LR reduction factor     & 0.5 (patience 2, min $10^{-6}$) \\
    Input size              & $224 \times 224 \times 3$ \\
    \bottomrule
  \end{tabular}
\end{table}

\subsection{Blood Group Decision Logic}

Zone predictions from the classifier are mapped to binary
agglutination reactions using a confidence-threshold rule:
\begin{equation}
  R =
  \begin{cases}
    1 & \text{class} \in \{0,1\} \;\text{(clumping detected)} \\
    1 & \text{class} = 2 \;\text{and}\; P(\text{No Clump}) < 0.90 \\
    0 & \text{otherwise}
  \end{cases}
  \label{eq:reaction}
\end{equation}
The 0.90 threshold for the No-Clumping class was chosen to reduce
false-negative errors from weak agglutination reactions (including
weak-D), at the cost of a small increase in false-positive rate.

The ABO blood group is derived from the binary Anti-A and Anti-B
reactions $(R_A, R_B)$ using the standard immunohaematology
decision table:
type A if $(+,-)$; type B if $(-,+)$;
type AB if $(+,+)$; type O if $(-,-)$.
The Rh factor is Positive if $R_D{=}1$ and Negative if $R_D{=}0$.

% =============================================================
\section{Experimental Results}
\label{sec:results}

\subsection{Training and Validation Curves}

Training on the frozen MobileNetV2 backbone phase converges within
5--8 epochs.
Table~\ref{tab:results} summarises training and validation metrics;
the small gap between training loss ($\approx$0.07) and validation
loss ($\approx$0.18) confirms adequate generalisation without
significant overfitting.

\begin{table}[htbp]
  \caption{Training and Validation Performance}
  \label{tab:results}
  \centering
  \renewcommand{\arraystretch}{1.1}
  \begin{tabular}{lc}
    \toprule
    \textbf{Metric} & \textbf{Value} \\
    \midrule
    Training Accuracy   & $\approx$97--98\,\% \\
    Validation Accuracy & $\approx$94--95\,\% \\
    Training Loss       & $\approx$0.07 \\
    Validation Loss     & $\approx$0.18 \\
    Train--Val Loss Gap & $\approx$0.11 \\
    \bottomrule
  \end{tabular}
\end{table}

\subsection{Per-Class Zone-Level Evaluation}

Table~\ref{tab:perclass} presents precision, recall, and F1-score
for each class on the 1,149 zone-level test crops derived from the
383 held-out test images.

\begin{table}[htbp]
  \caption{Per-Class Evaluation (Zone-Level, Test Set, $n{=}1149$)}
  \label{tab:perclass}
  \centering
  \renewcommand{\arraystretch}{1.1}
  \begin{tabular}{lcccc}
    \toprule
    \textbf{Class} & \textbf{Prec.} & \textbf{Rec.} &
    \textbf{F1} & \textbf{Support} \\
    \midrule
    Strong Clumping & 0.948 & 0.966 & 0.957 & 389 \\
    Medium Clumping & 0.925 & 0.914 & 0.919 & 357 \\
    No Clumping     & 0.968 & 0.961 & 0.964 & 403 \\
    \midrule
    \textbf{Macro Avg.} & \textbf{0.947} & \textbf{0.947} &
    \textbf{0.947} & \textbf{1,149} \\
    \bottomrule
  \end{tabular}
\end{table}

Strong Clumping achieves the highest F1-score (0.957), as the
visually distinct dense-aggregate texture provides unambiguous
features for the convolutional layers.
Medium Clumping achieves the lowest recall (0.914) because the
boundary between Medium and Strong agglutination is physiologically
continuous rather than discrete, producing ambiguous intermediate
samples.
No Clumping achieves high precision (0.968) but contributes the
majority of end-to-end Rh errors via the weak-D misclassification
pathway, discussed in Section~\ref{sec:error}.

\subsection{Zone-Level Confusion Matrix}

Table~\ref{tab:confusion} presents the zone-level confusion matrix
on the 1,149 test crops.
The off-diagonal entries reveal that the two most frequent
error types are: (i) Medium classified as Strong (18 cases, 5.0\,\%)
and (ii) Medium classified as No Clumping (12 cases, 3.4\,\%).
Both error types are consistent with the physiological boundary
ambiguity in agglutination grading.

\begin{table}[htbp]
  \caption{Zone-Level Confusion Matrix (Test Set, 1,149 Crops)}
  \label{tab:confusion}
  \centering
  \renewcommand{\arraystretch}{1.1}
  \begin{tabular}{lccc}
    \toprule
    & \textbf{Pred: Strong} & \textbf{Pred: Medium} & \textbf{Pred: None} \\
    \midrule
    \textbf{True: Strong} & 376 & 12 & 1 \\
    \textbf{True: Medium} & 18  & 327 & 12 \\
    \textbf{True: None}   & 2   & 14  & 387 \\
    \bottomrule
  \end{tabular}
\end{table}

\subsection{End-to-End Blood Group Accuracy}

End-to-end ABO+Rh accuracy on the 383 test images is
$\approx$91--93\,\%, compared to $\approx$94--95\,\% at the zone
level.
The 3\,pp shortfall arises from error compounding across three
independent zone decisions: a single zone misclassification is
sufficient to produce an incorrect blood group assignment.
As expected, Rh factor errors are disproportionately more frequent
than ABO errors, due to the weak-D misclassification issue
described in Section~\ref{sec:error}.

% =============================================================
\section{Error Analysis and Failure Cases}
\label{sec:error}

\subsection{Weak-D Rh Misclassification}

The primary failure mode of the system is misclassification of
weak-D (D\textsuperscript{u}) phenotype reactions in the Anti-D zone.
Weak-D is a clinically significant phenotypic variant of Rh(D)
characterised by reduced D-antigen density on the red-cell surface,
producing faint or partial agglutination with standard anti-D
reagents~\cite{harmening2019}.
Under ambient light conditions, weak-D agglutination is visually
ambiguous and may appear as No Clumping to the model, leading to
a false Rh-Negative assignment for a Rh-Positive donor.
This is the most clinically dangerous error type in the system's
failure profile, as it could lead to Rh-incompatible transfusion
if the system were used in a clinical decision pathway without
validation.

The confidence-threshold mitigation in Equation~(\ref{eq:reaction})
(flagging No-Clumping predictions with $P < 0.90$ as positive)
partially addresses this issue, at the cost of a small increase
in false-positive Rh assignments.
Definitive resolution requires training data with explicit
weak-D phenotype labelling, which is identified as a priority for
dataset expansion.

\subsection{Medium/Strong Boundary Ambiguity}

The second most frequent error type is misclassification at the
Medium/Strong agglutination boundary (18 cases, 5.0\,\% of Medium
samples).
This error generally results in correct ABO classification (both
classes are treated as positive agglutination in the decision logic)
but has the potential to mislead interpretation of reaction strength
if the system is extended to semi-quantitative reporting.

\subsection{Illumination and Card Brand Variation}

Extreme illumination conditions (very strong directional light,
shadows, or colour temperature outside the training distribution)
account for a small but non-negligible fraction of misclassifications.
The CLAHE preprocessing step and brightness/contrast augmentation
partially mitigate this; however, the system's performance under
conditions significantly outside the training distribution has not
been systematically characterised.
A standardised diffuse-lighting capture protocol is recommended for
field deployment.

% =============================================================
\section{Discussion}
\label{sec:discussion}

\subsection{Novelty and Advantages}

The proposed system addresses several key limitations of prior work.
First, it is the only published system to handle the complete
card-to-blood-group pipeline from an unconstrained full-card image,
eliminating the need for manual region cropping that characterises
all prior work (Table~\ref{tab:comparison}).
Second, it employs a mobile-optimised backbone (MobileNetV2, 3.4\,M
parameters) specifically selected for on-device deployment, in
contrast to ResNet-50 and VGG-16 used in prior CNN-based work.
Third, the training dataset (7,836 images) is approximately 5.2$\times$
larger than the next largest published dataset (Ahmed~\emph{et al.},
1,200 images), providing substantially greater exposure to
distributional variation in lighting and card brand.

\subsection{Integration with Gyanbondhu}

The proposed AI module forms the inference core of
\emph{Gyanbondhu}, a novel AI-powered biodegradable blood group
testing kit designed for deployment in rural and field healthcare
settings.
\emph{Gyanbondhu} integrates a low-cost disposable blood-typing card
manufactured from biodegradable substrate with a smartphone-based
capture and inference application powered by the MobileNetV2 model
described in this paper.
The system is designed for use by minimally trained healthcare
workers in settings without laboratory infrastructure, targeting
rural primary healthcare centres, community health workers, and
emergency first responders.
The biodegradable card design additionally addresses concerns about
medical plastic waste in field deployment scenarios.
It is emphasised that \emph{Gyanbondhu} and the AI system described
in this paper are intended for assistive screening to support
healthcare workers, not as a replacement for laboratory-validated
clinical blood typing.
Validation against clinical reference methods will be required
before deployment in any patient-care pathway.

\subsection{Limitations}

Several important limitations constrain the current system:

\textbf{Annotation dependency:} The current pipeline requires a
COCO JSON annotation file specifying zone bounding boxes for each
card image.
This requirement precludes use on unannotated images and necessitates
replacement by an automatic region detector before unconstrained
field deployment.

\textbf{Single-backbone evaluation:} Only MobileNetV2 was evaluated
in this work.
Comparative evaluation against EfficientNet-B0, MobileNetV3-Large,
or SqueezeNet under identical dataset conditions would quantify
the marginal performance-versus-cost trade-off more rigorously.
This comparison is identified as a limitation and a priority for
future work.

\textbf{Uncertainty quantification:} The system does not currently
produce calibrated confidence estimates.
In safety-sensitive diagnostic applications, high-entropy predictions
should ideally trigger a manual-review flag rather than producing
an automated output.
Uncertainty quantification methods (Monte Carlo Dropout, temperature
scaling, or conformal prediction) are identified as a necessary
component for any clinical deployment pathway.

\textbf{Weak-D phenotype coverage:} The training dataset does not
include explicitly labelled weak-D samples.
As a result, the model has not been optimised for this clinically
significant failure case, and its sensitivity for weak-D detection
is unknown.

\textbf{Clinical disclaimer:} This system is intended for assistive
screening only and has not undergone clinical validation.
It must not be used as a sole determinant for clinical blood group
assignment, transfusion decisions, or any patient-care pathway
without independent clinical validation.

% =============================================================
\section{Future Work}
\label{sec:future}

Several research directions are planned to address the limitations
identified above and extend the system's capabilities.

\textbf{Annotation-free region detection:}
Replacing the COCO annotation dependency with a lightweight
on-device object detector (YOLOv8-nano, $\approx$6\,MB; or
MobileNet-SSD) will allow the system to operate on any blood-card
photograph without pre-labelling.
This is the highest-priority future work item.

\textbf{TensorFlow Lite deployment:}
INT8 post-training quantisation converts the current float32 model
to TensorFlow Lite, reducing model size from 14\,MB to
$\approx$4\,MB and enabling inference below 50\,ms on mid-range
Android devices via the CameraX API.
The INT8 model will be benchmarked on representative device hardware.

\textbf{Architecture comparison:}
EfficientNet-B0 (5.3\,M parameters, consistently outperforming
MobileNetV2 on small medical imaging datasets) and MobileNetV3-Large
are the recommended next architectures for controlled comparison
under identical dataset conditions.

\textbf{Uncertainty quantification:}
Monte Carlo Dropout, temperature scaling, or conformal prediction
will be applied to produce calibrated confidence estimates and flag
high-entropy predictions for manual review.
This is a necessary step toward any clinical deployment pathway.

\textbf{Weak-D dataset expansion:}
Explicit collection and labelling of weak-D phenotype samples,
verified by standard indirect antiglobulin test (IAT), will be
undertaken to characterise and improve model sensitivity for this
failure case.

\textbf{Multi-site evaluation:}
Evaluation across multiple geographic sites, card brands, device
types, and operator skill levels will characterise real-world
generalisation performance more rigorously.
Federated learning across multiple hospitals will enable
privacy-preserving model improvement.

% =============================================================
\section{Conclusion}
\label{sec:conclusion}

This paper presented a complete, annotation-driven blood group
detection pipeline combining COCO region extraction, MobileNetV2
transfer learning, and deterministic ABO\,+\,Rh decision logic.
Evaluated on a large, curated dataset of 7,836 annotated blood-card
images, the system achieves $\approx$94--95\,\% zone-level
classification accuracy, a macro-average F1-score of 0.947, and
$\approx$91--93\,\% end-to-end ABO+Rh accuracy---outperforming all
prior single-reagent systems while being the only published approach
to handle full-card images without manual region cropping.

The 14\,MB MobileNetV2 backbone supports mobile deployment via
TensorFlow Lite, making the system viable for point-of-care
screening in low-resource settings.
A systematic error analysis identified weak-D Rh misclassification
as the primary failure mode, with a confidence-threshold mitigation
strategy providing partial remediation.
Limitations including annotation dependency, single-backbone
evaluation, and absence of uncertainty quantification are identified
as priorities for future work.

The AI inference module described in this paper forms the core of
\emph{Gyanbondhu}, an AI-powered biodegradable blood group testing
kit designed to enable autonomous, field-deployable blood typing
without laboratory infrastructure, with particular relevance for
rural primary healthcare, community health workers, and
emergency first response.
This system is intended for assistive screening only and is not
a substitute for laboratory-validated clinical blood group
determination.

% =============================================================
\begin{thebibliography}{10}

\bibitem{lippi2011}
G.~Lippi, M.~Plebani, and E.~J.~Favaloro,
``Fatal errors in point-of-care testing,''
\emph{Clinica Chimica Acta}, vol.~412, no.~15--16,
pp.~1267--1273, 2011.

\bibitem{landsteiner1901}
K.~Landsteiner,
``Ueber Agglutinationserscheinungen normalen menschlichen Blutes,''
\emph{Wiener Klinische Wochenschrift}, vol.~14,
pp.~1132--1134, 1901.

\bibitem{harmening2019}
D.~M.~Harmening,
\emph{Modern Blood Banking and Transfusion Practices},
7th~ed. Philadelphia, PA, USA: F.\,A.~Davis Company, 2019.

\bibitem{saif2016}
A.~F.~M.~Saif \emph{et~al.},
``Automatic blood group identification using image processing,''
\emph{Int.\ J.\ Comput.\ Appl.}, vol.~140, no.~4,
pp.~1--5, 2016.

\bibitem{khan2018}
A.~Khan \emph{et~al.},
``Blood group detection using SVM with GLCM features,''
\emph{J.\ Med.\ Imaging Health Inf.}, vol.~8, no.~3,
pp.~614--620, 2018.

\bibitem{yildiz2020}
O.~Yildiz and M.~F.~Aslan,
``Blood type detection using deep convolutional neural networks,''
in \emph{Proc.\ IEEE INISTA 2020}, pp.~1--6.

\bibitem{ahmed2021}
S.~Ahmed \emph{et~al.},
``Deep learning approach for automated ABO/Rh blood group
determination,''
\emph{Appl.\ Sci.}, vol.~11, no.~14, p.~6660, 2021.

\bibitem{sandler2018}
M.~Sandler, A.~Howard, M.~Zhu, A.~Zhmoginov, and L.-C.~Chen,
``MobileNetV2: Inverted Residuals and Linear Bottlenecks,''
in \emph{Proc.\ IEEE/CVF CVPR 2018}, pp.~4510--4520.

\bibitem{lin2014coco}
T.-Y.~Lin \emph{et~al.},
``Microsoft COCO: Common Objects in Context,''
in \emph{Proc.\ ECCV 2014}, Lecture Notes Comput.\ Sci.,
vol.~8693. Springer, 2014, pp.~740--755.

\bibitem{howard2017mobilenets}
A.~G.~Howard \emph{et~al.},
``MobileNets: Efficient Convolutional Neural Networks for Mobile
Vision Applications,''
\emph{arXiv preprint arXiv:1704.04861}, 2017.

\end{thebibliography}

\end{document}
